\section{Tolman--Eichenbaum Machine：What、Where 与 Associative Memory}

TEM 是第二讲的模型核心。它把 lateral entorhinal cortex（LEC）类比为 content/what stream，把 medial entorhinal cortex（MEC）类比为 structural/location stream，把 hippocampus 类比为 conjunctive associative memory。结构状态按 action recurrently 更新，内容在新环境中快速绑定，memory 再支持 recall 和 inference。

\subsection{架构总览：三种状态各司其职}

标题页之后，架构图展示三层信息：sensory observation、可复用 structural state 和 conjunctive memory。下方 graph 说明同一结构可承载不同 observation labels。模型不是把 observation 直接塞进一个 RNN，而是让 structure dynamics 与 content encoding 分别学习，再在 hippocampal-like state 中结合。

\lecturefigure{slide-50-tem-title.jpg}{Tolman--Eichenbaum Machine：统一空间与关系记忆}{官方录像手写教学状态 V050；00:59:38}

\lecturefigure{slide-51-tem-architecture.jpg}{TEM 架构：sensory content、structural state 与 memory}{官方录像手写教学状态 V051；01:01:08}

\begin{longtable}{@{}L{0.17\textwidth}L{0.31\textwidth}L{0.40\textwidth}@{}}
\toprule
组件 & 计算角色 & 神经科学类比 \\
\midrule
$x_t$ / sensory code & 当前 observation 的内容表示 & LEC-like what/content stream \\
$g_t$ / structural code & 由动作与上一状态更新的位置/关系表示 & MEC-like where/structure stream \\
$p_t$ / conjunctive code & 把当前内容与结构绑定为 episode & hippocampal conjunctive representation \\
$M_t$ / associative memory & 快速写入并按 cue 召回 conjunctive states & hippocampal recurrent/associative memory \\
\bottomrule
\end{longtable}

\subsection{LEC 与 MEC：内容和结构分离}

LEC 页从 observation 提取 sensory code，回答“看到了什么”；MEC 页根据 action/transition 更新 structural code，回答“在结构中的哪里”。这种分离让模型在新环境中保留 transition dynamics，同时替换 object identity。若二者从一开始完全纠缠，每次换 observation 都可能破坏已有地图。

\lecturefigure{slide-52-lec-what-stream.jpg}{LEC-like stream：编码 sensory content}{官方录像手写教学状态 V052；01:01:20}

\lecturefigure{slide-53-mec-where-stream.jpg}{MEC-like stream：更新 reusable structural state}{官方录像手写教学状态 V053；01:02:24}

结构更新可抽象为
$$
g_t=f_{\theta}(g_{t-1},a_t),
\qquad
x_t=e_{\phi}(o_t),
$$
其中 $o_t$ 是 observation，$a_t$ 是 action/edge，$g_t$ 是 path-integrated structural state，$x_t$ 是 content encoding。$\theta$ 跨环境慢速学习，当前环境中的 object binding 则可快速建立。

\begin{warningbox}{“Where” 不一定是二维坐标}
在 tree 或 transitive-order task 中，$g_t$ 表示的是 relational position，不必对应欧氏 $(x,y)$。把 MEC-like state 解释成通用结构码，正是 TEM 从 navigation 推广到 relational memory 的关键；但该推广仍需任务和神经数据支持。
\end{warningbox}

\subsection{Hippocampal binding 与 Hebbian-style memory}

Memory-update 页把 content 与 structure 结合成 conjunctive code，随后写入 associative memory。Hebbian 页给出“共同激活则连接增强”的局部更新直觉。一个抽象写法是
$$
p_t=b(x_t,g_t),
\qquad
M_t=\lambda M_{t-1}+p_t p_t^{\top},
$$
其中 $b$ 是 binding operation，$p_t$ 是 conjunctive representation，$M_t$ 是 memory matrix，$\lambda$ 控制旧记忆保留。真实 TEM 使用更细的 inference/generative 结构，但这个简式揭示快速写入为何不必重训全部参数。

\lecturefigure{slide-54-recurrent-memory-update.jpg}{TEM 的 recurrent state 与 associative-memory update}{官方录像手写教学状态 V054；01:02:58}

\lecturefigure{slide-55-hebbian-memory-equations.jpg}{Hebbian-style binding 与 memory equations}{官方录像手写教学状态 V055；01:03:40}

\begin{knowledgebox}{什么是 fast weights}
模型参数 $\theta,\phi$ 通过许多环境慢速训练，memory matrix $M_t$ 则在当前 episode 中快速更新，因此常被称为 fast weights。它允许一次或少量 exposure 后保存新内容，而结构 dynamics 仍由长期参数提供。
\end{knowledgebox}

\subsection{Graph inference 与 planning}

Graph-inference 页展示 agent 如何利用 structural state 和 recalled memory 推断未直接观察的节点或规划路径。关键不是把完整 graph 存成固定 adjacency table，而是让 recurrent transition model 生成候选 structure，再用 memory 补回当前环境的 observation identity。

\lecturefigure{slide-56-graph-inference-plan.jpg}{利用 structural state 与 memory 进行 graph inference}{官方录像手写教学状态 V056；01:06:02}

\begin{lstlisting}[caption={结构状态与快速绑定的概念流程}]
def tem_step(structure, action, observation, memory):
    next_structure = transition_model(structure, action)
    content = content_encoder(observation)
    conjunctive = bind(next_structure, content)
    memory.write(conjunctive)
    recalled = memory.read(next_structure)
    return next_structure, recalled
\end{lstlisting}

\teachervoice{Will 把目标描述成“结构慢学、内容快绑”。这与前半 SDM 的 address/value 分离形成呼应：structure 像可迁移 address system，sensory content 像要写入的 value，hippocampal memory 则把二者临时结合。}

\subsection{Grid-like code 与快速适应}

Grid-code 两页比较模型 representation 与神经数据中常见的 periodic spatial pattern；learning curve 页显示 agent 在少量访问后迅速改善。正确结论是结构化 recurrent state 与任务训练可以产生 grid-like basis，并支持快速 map learning。它不意味着每个 hidden unit 都对应真实 grid cell，也不证明六角格是唯一可行编码。

\lecturefigure{slide-57-grid-code-evidence.jpg}{TEM 中出现的 grid-like representations}{官方录像手写教学状态 V057；01:08:26}

\lecturefigure{slide-58-grid-code-summary.jpg}{模型与 neural grid codes 的结构比较}{官方录像手写教学状态 V058；01:09:36}

\lecturefigure{slide-59-quick-learning-result.jpg}{在少量节点访问后快速提升的学习曲线}{官方录像手写教学状态 V059；01:10:48}

\begin{knowledgebox}{读图：grid-like 不只看“像不像六边形”}
应检查 firing field 的周期性、方向、尺度、相位与环境变化下的稳定性，并比较 population structure。单张热图出现重复斑点可能来自平滑、边界或训练数据；只有多项几何统计与任务因果一致，才是更强证据。
\end{knowledgebox}

\subsection{本章小结}

TEM 用 content stream、structural recurrent state 与 fast associative memory 实现快速地图学习。结构表示跨环境复用，新的 observation 被快速绑定，memory 支持 graph inference。Grid-like code 和 few-visit learning 是功能证据，但不能被写成解剖身份。

